\section{Background and Measurements}\label{sec:measure}

\paragraph{Speculation and its break-even point.} A drafter proposes $k$ tokens and the target verifies them in one step, keeping the longest accepted prefix plus one token of its own. We use EAGLE-3 drafters \citep{li2025eagle3} with $k{=}3$, so the acceptance length $a$, the number of tokens emitted per verification step, lies between 1 and 4. At concurrency $C$ (requests being decoded), speculation is profitable when its throughput exceeds that of plain decoding, i.e., when the ratio $G(a,C)$ of the two exceeds one. Verification makes every step more expensive, and at high concurrency the GPU has no spare compute to absorb this cost, so the break-even acceptance $a^{*}(C)$ at which $G=1$ rises with load.

\paragraph{Measurement setup.} The measurements in this section use vLLM 0.27.1 on one H100 GPU, greedy decoding, the factory EAGLE-3 drafters for Qwen3-8B and Qwen3-30B-A3B released with AngelSlim \citep{yang2025qwen3,angelslim}, five short-output domains (chat, SQL, code, Chinese, function calling) and MATH-500 with thinking for long outputs (\cref{app:workloads}). They are the calibration inputs of the controller and were taken before the evaluation checkpoint was frozen (\cref{sec:setup}).

\paragraph{Whether speculation pays depends on load and architecture.} Figure~\ref{fig:motivation}a sweeps concurrency. For Qwen3-8B with short outputs, speculation is \nv{sweep.short.c1}\% faster than plain decoding at $C{=}1$ but \nv{sweep.short.c256abs}\% slower at $C{=}256$; with long outputs it remains \nv{sweep.long.c96}\% faster at the largest load we could run ($C{=}96$). The MoE target behaves the other way: the gain grows from \nv{sweep.moe.c1}\% at $C{=}1$ to \nv{sweep.moe.c64}\% at $C{=}64$ and is still \nv{sweep.moe.c256}\% at $C{=}256$. At $C{=}256$ a speculative step costs \nv{step.m2.c256}$\times$ a plain step on the MoE target but \nv{step.m1.c256}$\times$ on the dense one, consistent with the observation that sparse MoE targets profit from speculation at larger batches \citep{huang2025moesd}. A rule such as ``disable speculation under heavy load'' is right for one target and badly wrong for the other.

\paragraph{Co-located training is not free.} The trainer shares the GPU by time slicing: at duty cycle $d$ it computes for a fraction $d$ of the time. Serving throughput falls almost linearly, by $\kappa(C)\,d$ with $\kappa$ between \nv{kappa.min} and \nv{kappa.max} (Figure~\ref{fig:motivation}b): duty 0.1 costs \nv{duty.d01.lossmin}--\nv{duty.d01.lossmax}\% of throughput, duty 0.5 costs \nv{duty.d05.lossmin}--\nv{duty.d05.lossmax}\%, and training all the time costs \nv{duty.d10.lossmin}--\nv{duty.d10.lossmax}\% and raises p99 time per output token (TPOT) by \nv{duty.d10.p99min}--\nv{duty.d10.p99max}\%.

\paragraph{Learning ceilings differ.} Figure~\ref{fig:motivation}c shows the acceptance length after learning online from $N$ requests, one pass in arrival order, with the target's outputs as labels. Code and SQL improve steadily. Chinese stalls at \nv{lc.zh.n3200}, because only \nv{cover.zh}\% of the target's output tokens are in the drafter's 32k vocabulary. Function calling peaks at $N{=}\nv{lc.fc.peakN}$ and then declines. Against the break-even line at $C{=}256$, the code drafter crosses after about 100 requests, the Chinese drafter never does, and the MoE drafter is above its much lower break-even (\nv{astar.moe.c256}) from the start.

\paragraph{The profitability frontier.} For each domain $\ell$ we measured the throughput ratio of three or four drafter checkpoints at six loads. At every load the ratio is close to linear in acceptance (\cref{app:calibration}, Figure~\ref{fig:frontier}), so we fit
\begin{equation}
\Ghat_\ell(a,C)=\alpha_\ell(C)\,a+\beta_\ell(C),\qquad a^{*}_\ell(C)=\frac{1-\beta_\ell(C)}{\alpha_\ell(C)},\label{eq:frontier}
\end{equation}
with residuals of at most \nv{astar.code.err} for code, \nv{astar.zh.err} for Chinese, and \nv{astar.moe.err} for MoE code (\nv{astar.sql.err} at one SQL load). For Qwen3-8B the break-even acceptance rises from 1.0--1.2 at low load to \nv{astar.code.c256} (code), \nv{astar.sql.c256} (SQL), and \nv{astar.zh.c256} (Chinese) at $C{=}256$; for the MoE target it is \nv{astar.moe.c256}. The adapted code drafter reaches \nv{fr.code.n3200.c256}$\times$ plain decoding at $C{=}256$ in this sweep, whereas the factory drafter reaches \nv{fr.code.n0.c256}$\times$.
